% 2.3 — Les bacs à sable
\section{Les bacs à sable}
\label{sec:sandbox}

Un bac à sable confie au système d'exploitation le soin de limiter ce qu'une commande peut
atteindre. Les deux agents en proposent un.

\begin{description}[style=nextline,leftmargin=0.6cm,font=\sffamily\bfseries]
  \item[Claude Code]
    Le bac à sable de l'outil shell s'appuie sur Seatbelt sous macOS et sur
    \cmd{bubblewrap} sous Linux et WSL2. Il est désactivé par défaut. Une fois activé, les
    commandes peuvent écrire dans le répertoire de travail, un répertoire temporaire et les
    répertoires ajoutés ; certains chemins sensibles restent protégés en écriture (fichiers
    de configuration de l'agent, \fichier{.bashrc}, \textit{hooks} git\ldots). Sous Windows
    natif, les commandes s'exécutent sans bac à sable~\cite{cc-sandbox}.
  \item[Codex]
    Codex sépare deux réglages : une politique d'approbation (quand demander) et un mode de
    bac à sable (ce que le système autorise), dont le mode \cmd{workspace-write} qui limite
    les écritures à l'espace de travail~\cite{codex-security}. Sous Linux, l'isolation du
    système de fichiers repose par défaut sur \cmd{bubblewrap}~\cite{codex-linux-sandbox}.
  \item[Conteneurs et machines virtuelles]
    Faire travailler l'agent dans un conteneur ou une machine virtuelle jetable isole la
    machine hôte entière.
\end{description}

\paragraph{Ce qu'un bac à sable ne couvre pas.} Le bac à sable répond à la question
« \emph{où} la commande peut-elle écrire ? ». Or le projet lui-même --- le répertoire de
travail --- est, par construction, inscriptible : c'est là que l'agent travaille. Un
\cmd{rm -rf src} ou un \cmd{mv} qui écrase un fichier non suivi par git y est autorisé.
De même, une commande qui atteint une base de données de production avec des identifiants
valides, ou qui pousse de force sur un dépôt distant, utilise un accès que le bac à sable
a laissé ouvert pour que l'agent puisse travailler. Enfin, la protection dépend de la
plateforme : là où le bac à sable n'existe pas (Windows natif pour Claude Code), il n'y a
rien.

\begin{pourquoi}[Complémentaire, pas concurrent]
Le bac à sable \emph{borne} ; \demo{} \emph{juge et récupère à l'intérieur} de la borne.
Les deux se combinent : le bac à sable empêche l'agent de sortir du projet, \demo{}
garantit qu'à l'intérieur du projet une destruction reste réversible, et qu'une action à
distance passe par une personne.
\end{pourquoi}
